\begin{qfdefinition}[Crude oil markets: grades, WTI and Brent, OPEC and spare capacity, chokepoints and pipelines, and the 3:2:1 crack spread, as of 2012--2013]\label{def:crude-oil-markets}
\qffield{Statement} Oil supplied 33 percent of world energy demand in 2012, coal 30 and natural gas 24 \citep[chap.~10, sec.\ ``Energy'', PDF p.~301]{s7-gliner}. Quality is measured by API gravity (higher means lighter; below 10 it sinks in water; 25 or less is ``heavy'') and sulphur content (sweet against sour), light sweet grades being the most valuable (pp.~301--302). OPEC has 12 members producing over 35 percent of world output and holding more than 70 percent of reserves; Saudi Arabia has the largest spare capacity, so ``the Call on OPEC'' usually means Saudi production rises (pp.~302--305, Table~10.4, OPEC-12 quota 30 million barrels a day as of September 2013).

WTI, the U.S.\ benchmark, is delivered at Cushing, Oklahoma, and trades on NYMEX (CL), 1,000 barrels a contract; Brent, from North Sea fields, ``is used to price more than half of the world's internationally traded crude oil'' and trades on ICE (CO), 1,000 barrels; Brent's premium over WTI in Gliner's time reflected North Sea depletion, seaborne supply risk and pipeline constraints at Cushing (pp.~305--307). World consumption was 90 million barrels a day in 2012; the Strait of Hormuz carried 17 million barrels a day in 2011, ``about 20 percent of global traded oil and 35 percent of seaborne oil'', Malacca 14 million, Bab el-Mandeb 3 million and the Suez Canal 2 million; the Druzhba pipeline runs over 2,300 miles and carries up to 1.4 million barrels a day (pp.~308--310).

The \emph{crack spread} is the refining margin; the 3:2:1 spread (three barrels of crude against two of gasoline and one of heating oil) is the most popular because gasoline output is almost double heating oil's, with product prices in cents a gallon converted at 42 gallons a barrel (sec.\ ``Crack Spread'', pp.~311--312).

\qffield{Notation} 3:2:1 crack per barrel $=\big[2\times(\text{RBOB}\times0.42)+1\times(\text{HO}\times0.42)-3\times\text{CL}\big]/3$, products in cents per gallon and crude in dollars per barrel; Gliner's worked example gives \$27.95 a barrel (p.~312), and crude at \$80, gasoline at \$2.40 and heating oil at \$2.70 a gallon give \$25.00.

\qffield{Conditions} All shares, capacities, membership and contract details are Gliner's, as of 2011--2013. Benchmark composition and trade flows may have changed since. \emph{Membership}. OPEC's home page lists twelve Member Countries in its navigation: Algeria, Congo, Equatorial Guinea, Gabon, Iran, Iraq, Kuwait, Libya, Nigeria, Saudi Arabia, the United Arab Emirates and Venezuela \citep[home page navigation]{opec-home-2026}.

\emph{The UAE}. Four kinds of evidence, kept apart. (i) The announcement: the UAE's state news agency reported on 28 April 2026 that the UAE ``announced its decision to exit the Organisation of the Petroleum Exporting Countries (OPEC and OPEC+), effective 1 May 2026'', having joined in 1967 through Abu Dhabi \citep[pp.~1--2]{wam-uae-opec-2026}.

(ii) EIA's classification: its analysis of 23 June 2026 treats the UAE as having left effective 1 May, and puts OPEC's 2025 crude output, UAE included, at 28.0 million barrels a day, 35 percent of the world total, or 31 percent without the UAE, and OPEC+ at about 46 percent, about 42 without it \citep[web page]{eia-uae-opec-2026}; its live glossary lists eleven current members without the UAE \citep[glossary entry]{eia-opec-glossary}.

(iii) OPEC's own pages: the site's navigation, on the home page and on a press release of 3 May 2026, still lists the UAE among twelve Member Countries, and that release, on a June production adjustment of 188 thousand barrels a day by seven OPEC+ countries (Saudi Arabia, Russia, Iraq, Kuwait, Kazakhstan, Algeria and Oman), says nothing about membership \citep[release and site navigation]{opec-pr-2026-05-03}.

(iv) The statutory rule: under Article~8A of OPEC's Statute a withdrawal notice ``shall take effect at the beginning of the next calendar year after the date of its receipt by the Conference, subject to the Member having at that time fulfilled all financial obligations arising out of its membership'' \citep[2021 edition, Art.~8A, PDF p.~12]{opec-statute-2021}, which on its face would put the legal effect at 1 January 2027, not 1 May 2026, unless the Conference agreed otherwise.

No OPEC statement on the date or legal effect of the UAE's withdrawal was found, so whether the UAE is a member in law on 16 October 2026 is open; for market data, EIA already counts OPEC without it. EIA also notes the conflict in Iran from 28 February 2026 and the effective closure of the Strait of Hormuz, so Gliner's chokepoint flows above are historical.

\qffield{Practice} The crack spread is a margin before operating costs; refiners' actual yields and product slates differ from 3:2:1. Political risk in producers and chokepoint disruptions move Brent first, in Gliner's account.

\qffield{Pitfalls} \emph{Terminology.} Gliner writes that ``kerosene is used for diesel'' (p.~301); kerosene is jet fuel's base and diesel is a separate distillate. Gliner calls Saudi Aramco ``private'' and lists it among the ``biggest IOCs'' (pp.~303--304); it is a state-owned national oil company. Gliner describes Brent as ``deliverable to Sullom Voe'' (p.~306). ICE now describes its Brent Crude futures as ``a deliverable contract based on EFP delivery with an option to cash settle against the ICE Brent Index price for the last trading day'', 1,000 barrels, trading to the last business day of the second month preceding the contract month \citep[contract specifications]{ice-brent-spec}: delivery is by exchange for physical, not to a terminal, and there is a cash-settlement option.

\emph{NYMEX WTI, in part}. NYMEX's certification to the CFTC of 27 January 2023 (Submission 23-057) reprints Rules 200101 and 200104 of the Light Sweet Crude Oil futures (CL) for contract months from April 2023: deliverable domestic crude is a pipeline's common Domestic Sweet stream of WTI-type light sweet crude with, among other limits, sulphur at most 0.42 percent by weight, API gravity from 37 to 42 degrees and Reid vapour pressure below 9.5 psi (the test method moving to ASTM D-6377), delivered free on board at any pipeline or storage facility in Cushing, Oklahoma with access to Enterprise, Enbridge or Plains storage, each ``including any successor in such capacity'' \citep[Exhibits A--B]{nymex-23-057}.

\emph{Current chapter} (official PDF as of October 2026, undated). Chapter~200 now names ASTM D-6377 for vapour pressure and keeps the 23-057 grade and Cushing terms; the trading unit is 1,000 U.S. barrels (42,000 gallons), with a delivery tolerance of 2 percent either way except for book-outs and in-tank or in-line transfers; the minimum fluctuation is \$0.01 a barrel (\$10 a contract); and trading ends on the third business day before the 25th calendar day of the month preceding delivery, or, if the 25th is not a business day, the third business day before the last business day preceding the 25th \citep[Rules 200101 and 200102.B--F]{nymex-ch200-2026}. Gliner says only that WTI is the higher grade; that it historically traded at a premium for that reason is an interpretation.

\qffield{See also} \hyperref[def:natural-gas-market]{natural gas}; \hyperref[def:contango-backwardation-stocks]{contango and backwardation}; \hyperref[def:commodity-indices-risk-return]{commodity indices}.

\qffield{Further reading} Benchmarks, price reporting and the negative WTI settlement of April 2020.

\end{qfdefinition}
